\documentclass{article}
\usepackage[utf8]{inputenc}
\usepackage[T1]{fontenc}
\usepackage{geometry}
\geometry{
    a4paper,
    left=2.5cm,
    right=2.5cm,
    top=2.5cm,
    bottom=2.5cm
}
\usepackage{graphicx}
\usepackage{float}
\usepackage{hyperref}
\usepackage[english]{babel}
\usepackage{amsmath}
\usepackage{amssymb}
\usepackage{array}
\usepackage{booktabs}
\usepackage{caption}
\usepackage{subcaption}
\usepackage{enumitem}
\usepackage{fancyhdr}
\usepackage{titlesec}
\usepackage{url}
\usepackage{textcomp}
\usepackage{xcolor}
\usepackage{listings}
\usepackage{multirow}
\usepackage{longtable}
\usepackage{tabularx}
\usepackage{times}

% Set up page style
\pagestyle{fancy}
\fancyhf{}
\rfoot{\thepage}
\renewcommand{\headrulewidth}{0pt}

% Title formatting
\titleformat{\chapter}[display]{\normalfont\bfseries}{\chaptertitle\thechapter}{0pt}{\Huge}
\titlespacing*{\chapter}{0pt}{50pt}{40pt}

% Document info
\title{Template Document}
\author{Academic Support Agent}
\date{\today}

\begin{document}
\maketitle

\maketitle % N o remova essa linha!

\begin{abstract} % resumo em ingl s
  Escreva seu resumo em l ngua estrangeira (ingl s)...
\end{abstract}
     
\begin{resumo}
  Resumo do trabalho (portugu s)...  
\end{resumo}

\section{Introdu o}
Cen rios, motiva o, justificativa, problema de pesquisa, objetivo geral e espec ficos (taxonomia de bloom);

\section{Referencial Te rico}
Voltado ao objetivo geral (teoria por tr s do m todo), deve conter os assuntos bases da pesquisa, fazendo cita es indiretas e diretas curtas.

\begin{comment}
testando comentario
\end{comment}

\begin{description}
    \item[Exemplos de cita o indireta]    
\end{description}

Segundo \textcite{spinello2024}, o trabalho de conclus o deve ter cita es retiradas de artigos cient ficos encontrados nas bases de dados. 

O trabalho de conclus o deve ter cita es retiradas de artigos cient ficos encontrados nas bases de dados \cite{spinello2024}.

\begin{description}
    \item[Exemplos de cita o direta curta]
\end{description}

Segundo \textcite{spinello2024} \enquote{o trabalho de conclus o deve ter cita es retiradas de artigos cient ficos encontrados nas bases de dados}. Note que para colocar um texto entre aspas, usamos o comando \verb|\enquote{texto}|.

Ressalta-se que o \enquote{trabalho de conclus o deve ter cita es retiradas de artigos cient ficos encontrados nas bases de dados} \cite{spinello2024}.

No estudo comparativo apresentado em \textcite[p. 107]{rabello2010} ...

No trabalho de \textcite{pargaonkar2021} ...

Nos trabalhos de \textcites{badgujar2024}{pargaonkar2021} s o aplicadas t cnicas de ...

No artigo de \textcite{estevao2023} ...

\textsc{Cita o direta longa devem ser evitadas em artigos cient ficos!}

\section{Trabalhos Relacionados}

Trabalhos semelhantes aos objetivos espec ficos, sempre detalhando ao final da se o a
diferen a ao trabalho proposto (quantidade -- 5 trabalhos);

Neste item ser o apresentados os principais trabalhos que possuem uma rela o com o assunto definido neste estudo....

% Itens com marcadores
\begin{itemize}
     \item \textbf{T tulo do artigo 01 \cite{ogliari2019}}
     
     Primeiro par grafo indicar uma introdu o do assunto...
     
     No segundo: o que o estudo procurou analisar, qual o objetivo...
     
     No terceiro: o que foi desenvolvido, qual aplica o/experimento foi realizado...
     
      ltimo: em quais conclus es o trabalho chegou...
    \newline
    
    \item \textbf{T tulo do artigo 02 (Autor, ano)}
    
     Primeiro par grafo indicar uma introdu o do assunto...
     
     No segundo: o que o estudo procurou analisar, qual o objetivo...
     
     No terceiro: o que foi desenvolvido, qual aplica o/experimento foi realizado...
     
      ltimo: em quais conclus es o trabalho chegou...
    
\end{itemize}

\section{Materiais e M todos}
    Tecnologias, instrumentos e procedimentos que ser o usados no estudo. O Algoritmo \ref{alg:id_algo} se refere ao m todo de ordena o Bubblesort expresso em linguagem Python.

    % C digo-fonte formatado
    % Ver: https://en.wikibooks.org/wiki/LaTeX/Source_Code_Listings
    \begin{lstlisting}[  
        %float,
        language=Python, 
        frame=single, 
        numbers=left,
        caption={M todo de ordena o Bubblesort},
        label={alg:id_algo} % id para referenciar
        ]        
    def bubble_sort(alist):
        for i in range(len(alist)-1,0,-1):
            for j in range(i):
                if alist[i]>alist[i+1]:
                    temp = alist[i]
                    alist[i] = alist[i+1]
                    alist[i+1] = temp    
    \end{lstlisting}

    % Figura
    % https://pt.overleaf.com/learn/latex/Inserting_Images
    \begin{figure}[H]
        \centering
        \includegraphics[width=0.5\textwidth]{fig1.jpg} % altere a width (intervalo entre 0.0 e 1), caso a imagem n o fique na posi o correta
        \caption{Minha figura}
        \label{fig:id_figura}
    \end{figure}
    
    % Tabela
    % Editor online: https://www.latex-tables.com
     \begin{table}[ht]
        \centering
        \caption{Minha tabela}
        \label{tab:id_tabela}
        \begin{tblr}{
          colspec = {X[l] X[l]}, % Tipo X define c lulas com quebra autom tica de linha. Use: l=left, r=right, c=center ou j=justify para alinhamento dentro das c lulas
          width = \linewidth,
          hlines, % inclui borda horizontal
          vlines, % inclui borda vertical        
        }
        \textbf{cabe alho 1} & \textbf{cabe alho 2} \\
        {texto   esquerda} & {Existem muitas varia es das passagens do Lorem Ipsum dispon veis, mas a maior parte sofreu altera es de alguma forma, pela injec o de humor, ou de palavras aleat rias que nem sequer parecem suficientemente cred veis. Se vai usar uma passagem do Lorem Ipsum, deve ter a certeza que n o cont m nada de embara oso escondido no meio do texto.}
        \end{tblr}
    \end{table}

\section{Resultados e Discuss o}
Essa se o dever  ser escrita na segunda parte do trabalho, conhecida como TCC2.

\section{Considera es Finais}
Essa se o dever  ser escrita na segunda parte do trabalho, conhecida como TCC2.

% %%%%%%%%%%%%%%%%%%%%%%%%%%%%%%%%%%%%%%%%%%%%%%%%%%%%%%%%%%%%%%
% Se o de Refer ncias (gerada automaticamente)
\printbibliography  % N o remover esta linha
% %%%%%%%%%%%%%%%%%%%%%%%%%%%%%%%%%%%%%%%%%%%%%%%%%%%%%%%%%%%%%%
\end{document}